% SHARD-ID: AQM-19-CANONICAL-BOSON-FERMION-FIELD-COMPARISON
% SHARD-TITLE: Canonical Boson and Fermion Fields as Field-Label Tests
% SHARD-SUMMARY: Tests whether standard finite bosonic and fermionic fields fit the arithmetic-field label-space pattern.
% SHARD-SUMMARY: Proves that finite bosonic fields arise from real symplectic map spaces and Stone--von Neumann uniqueness.
% SHARD-SUMMARY: Separates fermionic CAR fields from the naive proposal of replacing the target by Grassmann numbers.
% SHARD-KEYWORDS: bosonic field, fermionic field, CCR, CAR, Weyl, Stone-von-Neumann, Grassmann, Fock

\section{Canonical Boson and Fermion Fields as Field-Label Tests}
\label{sec:canonical-boson-fermion-field-comparison}

\subsection{Status and Source Anchors}
This shard is a kinematic comparison.  It does not introduce dynamics,
Hamiltonians, continuum limits, locality cones, or equations of motion.  The
question is only whether the field-label construction from
Shard \(AQM\)-14 reproduces the standard finite-mode canonical fields.

The CCR and CAR ground truth is local; paths and hashes are registered in
\path{references/canonical_fields/SOURCES.md}.  The line anchors used here are:
Derezinski's \path{derezinski.tex}, lines 1078--1135 for Weyl CCR, 1197--1284
for Schrodinger/Stone--von Neumann, 1437--1524 for CAR, 1612--1696 for finite
CAR representations, 1803--1964 and 2098--2131 for Fock and creation operators,
2251--2301 for bosonic Fock CCR fields, and 2629--2662 for fermionic Fock CAR
fields; Bekka's \path{CCR-GeneralRings-v5.tex}, lines 220--231, 253--273,
287--350, and 603--633; and Keyl--Schlingemann's \path{fidelity_arXiv.tex},
lines 253--260, 270--355, 411--492, and 496--531.

\subsection{Lamport Dependency Skeleton}
\[
\begin{array}{rcl}
D1&:&X\text{ is a finite set of sites.}\\
D2&:&V_B=\mathbb R_q\oplus\mathbb R_p,\quad
       \omega_B((q,p),(q',p'))=pq'-p'q.\\
D3&:&E_B=\operatorname{Map}(X,V_B),\quad
       \Omega_B=\sum_{x\in X}\omega_B.\\
D4&:&W_B:E_B\to U(\mathcal H)\text{ satisfies Weyl-form CCR.}\\
L1&:&(E_B,\Omega_B)\text{ is a real symplectic vector space.}\\
L2&:&\text{Irreducible regular }W_B\text{ is Schrodinger by Stone--von Neumann.}\\
T1&:&D1\text{--}D4\text{ reproduce the finite standard bosonic field.}\\[1mm]
D5&:&Z=\operatorname{Map}(X,K)\text{ is a finite complex one-particle space.}\\
D6&:&Y_F=\operatorname{Re}(Z\oplus\overline Z),\quad
       \alpha((z,\bar z),(w,\bar w))=\operatorname{Re}\langle z,w\rangle.\\
D7&:&\Gamma_a(Z)=\bigoplus_k\wedge^k Z,\quad
       \phi(z,\bar z)=a^*(z)+a(z).\\
L3&:&a,a^*\text{ satisfy the finite exterior-algebra CAR.}\\
L4&:&\{\phi(y),\phi(y')\}=2\alpha(y,y').\\
T2&:&D5\text{--}D7\text{ reproduce the finite standard fermionic field.}\\
L5&:&\text{Grassmann target replacement is not the same construction as }T2.
\end{array}
\]

\subsection{Bosonic Field Labels D1--D4}
Let \(X\) be finite.  Write
\[
  \mathcal Q_X=\operatorname{Map}(X,\mathbb R)
\]
for the real configuration space.  Its algebraic dual is identified with
\(\operatorname{Map}(X,\mathbb R)\) by the finite sum
\[
  p(q)=\sum_{x\in X}p_xq_x.
\]
Thus
\[
  E_B=\operatorname{Map}(X,\mathbb R_q\oplus\mathbb R_p)
      \cong \mathcal Q_X\oplus\mathcal Q_X^*.
\]
For \(\phi=(q,p)\) and \(\psi=(q',p')\), define
\[
  \Omega_B(\phi,\psi)
  =
  \sum_{x\in X}\bigl(p_xq'_x-p'_xq_x\bigr).
\]
A Weyl-form CCR representation of \(E_B\) is a map
\[
  W_B:E_B\to U(\mathcal H)
\]
such that
\[
  W_B(\phi)W_B(\psi)
  =
  \exp\!\left(-\frac{i}{2}\Omega_B(\phi,\psi)\right)W_B(\phi+\psi).
\]
This is Derezinski's convention with \(\eta=p\) and \(q=q\).

\paragraph{Lemma \(L1\).}
\((E_B,\Omega_B)\) is a finite-dimensional real symplectic vector space.
\emph{Proof.}
Pointwise addition and scalar multiplication make \(E_B\) a real vector space.
The formula for \(\Omega_B\) is a finite sum of bilinear expressions, hence is
bilinear.  Also
\[
  \Omega_B(\phi,\phi)=\sum_x(p_xq_x-p_xq_x)=0,
\]
so it is alternating.  If \(\phi=(q,p)\ne0\), then either some \(q_{x_0}\ne0\)
or some \(p_{x_0}\ne0\).  In the first case take \(\psi\) supported at \(x_0\)
with \(p'_{x_0}=1\) and \(q'_{x_0}=0\); then
\(\Omega_B(\phi,\psi)=-q_{x_0}\ne0\).  In the second case take \(\psi\)
supported at \(x_0\) with \(q'_{x_0}=1\) and \(p'_{x_0}=0\); then
\(\Omega_B(\phi,\psi)=p_{x_0}\ne0\).  Thus the radical is zero.

\paragraph{Lemma \(L2\).}
Every irreducible regular representation of the Weyl-form CCR over
\((E_B,\Omega_B)\) is unitarily equivalent to the Schrodinger representation on
\[
  L^2(\mathcal Q_X)\cong L^2(\mathbb R^{|X|}).
\]
\emph{Proof.}
Choose an ordering of \(X\).  Then \(\mathcal Q_X\cong\mathbb R^N\) for
\(N=|X|\), and Lemma \(L1\) identifies \(E_B\) with
\(\mathcal Q_X^*\oplus\mathcal Q_X\) carrying the canonical symplectic form.
Derezinski's finite-dimensional Stone--von Neumann theorem states that a
regular CCR representation over such a finite symplectic space is a
Schrodinger representation up to multiplicity, and irreducibility is exactly
the one-copy case.  Bekka's source records the same local-field uniqueness
principle and the irreducibility of the Schrodinger pair.

\paragraph{Theorem \(T1\).}
For finite \(X\), the \(AQM\)-14 field-label recipe with real target \(V_B\)
recovers the standard finite bosonic canonical field.
\emph{Proof.}
The input is exactly \(D1\)--\(D3\): a finite physical set, a target vector
space with a scalar-valued alternating form, and a field-label space of maps.
Lemma \(L1\) checks that the summed pointwise form is symplectic.  Definition
\(D4\) is the Weyl quantization step.  By Lemma \(L2\), the irreducible regular
output is, up to unitary equivalence, the Schrodinger representation on
\(L^2(\mathbb R^{|X|})\).  Derezinski's bosonic Fock source gives the equivalent
finite-mode Fock realization by creation and annihilation operators.  Therefore
the finite standard bosonic field is a direct real analogue of the arithmetic
finite-field Weyl construction, with \(\mathbb F_q\) replaced by \(\mathbb R\).

\subsection{Fermionic Field Labels D5--D7}
Now let \(K\) be a finite-dimensional complex internal one-particle Hilbert
space and set
\[
  Z=\operatorname{Map}(X,K).
\]
Its inner product is the finite sum
\[
  \langle z,w\rangle_Z=\sum_{x\in X}\langle z_x,w_x\rangle_K.
\]
The real CAR label space is
\[
  Y_F=\operatorname{Re}(Z\oplus\overline Z)
      =\{(z,\bar z):z\in Z\},
\]
equipped with
\[
  \alpha((z,\bar z),(w,\bar w))=\operatorname{Re}\langle z,w\rangle_Z.
\]
The fermionic Hilbert space is the antisymmetric Fock space
\[
  \Gamma_a(Z)=\bigoplus_{k=0}^{\dim Z}\wedge^k Z.
\]
For \(z\in Z\), define creation by \(a^*(z)\Psi=z\wedge\Psi\).  Define
annihilation \(a(z)\) as the adjoint of \(a^*(z)\), equivalently on a
decomposable wedge by
\[
  a(z)(w_1\wedge\cdots\wedge w_k)
  =
  \sum_{j=1}^k(-1)^{j-1}\langle z,w_j\rangle\,
  w_1\wedge\cdots\widehat{w_j}\cdots\wedge w_k .
\]
The real fermion field operator is
\[
  \phi(z,\bar z)=a^*(z)+a(z).
\]

\paragraph{Lemma \(L3\).}
The exterior creation and annihilation operators satisfy
\[
  \{a^*(z),a^*(w)\}=0,\qquad
  \{a(z),a(w)\}=0,\qquad
  \{a(z),a^*(w)\}=\langle z,w\rangle I .
\]

\emph{Proof.}
The first identity is the antisymmetry of the wedge product:
\[
  z\wedge w\wedge\Psi+w\wedge z\wedge\Psi=0.
\]
The second follows from the displayed contraction formula, since deleting two
vectors in opposite orders gives opposite signs.  For the mixed relation, use
the contraction identity
\[
  a(z)(w\wedge\Psi)=\langle z,w\rangle\Psi-w\wedge a(z)\Psi .
\]
Therefore
\[
\begin{aligned}
  a(z)a^*(w)\Psi+a^*(w)a(z)\Psi
  &=a(z)(w\wedge\Psi)+w\wedge a(z)\Psi\\
  &=\langle z,w\rangle\Psi .
\end{aligned}
\]
This proves the three CAR identities.

\paragraph{Lemma \(L4\).}
For \(y=(z,\bar z)\) and \(y'=(w,\bar w)\),
\[
  \{\phi(y),\phi(y')\}=2\alpha(y,y')I .
\]

\emph{Proof.}
Expand the anticommutator and apply Lemma \(L3\):
\[
\begin{aligned}
  \{\phi(y),\phi(y')\}
  &=
  \{a^*(z)+a(z),a^*(w)+a(w)\}\\
  &=\langle z,w\rangle I+\langle w,z\rangle I\\
  &=2\operatorname{Re}\langle z,w\rangle I
   =2\alpha(y,y')I .
\end{aligned}
\]

\paragraph{Theorem \(T2\).}
For finite \(X\), the standard finite fermionic canonical field is the CAR
construction \(D5\)--\(D7\), not a Weyl-Heisenberg quantization of a
Grassmann-valued target.

\emph{Proof.}
Derezinski's CAR definition uses a real Hilbert label space with positive
scalar product and field operators satisfying an anticommutator relation.
Definitions \(D5\)--\(D7\) instantiate this with \(Z=\operatorname{Map}(X,K)\).
Lemma \(L4\) proves the CAR relation on the antisymmetric Fock space, and
Derezinski's Fock source identifies this as the fermionic Fock representation
of the CAR.  Thus the finite fermionic field arises from the same first step
"choose site fields", but the quantization functor changes: symmetric
symplectic/Weyl data give bosons, while positive-symmetric Clifford/CAR data
give fermions.

\subsection{Grassmann Target Test L5}

\paragraph{Lemma \(L5\).}
Replacing \(V_B\) by a Grassmann algebra is not a direct construction of the
standard finite fermion field in the \(AQM\)-14 sense.

\emph{Proof.}
The \(AQM\)-14 Weyl construction requires an ordinary vector-space label
group, a scalar-valued alternating form, and a projective unitary
representation whose cocycle is a complex phase or finite-field character.
The standard fermion field in \(T2\) instead satisfies the CAR relation
\(\{\phi(y),\phi(y')\}=2\alpha(y,y')I\), with a symmetric positive scalar
product.  This is already a different algebraic input.

Keyl--Schlingemann show that Grassmann variables do enter fermionic phase-space
methods inside GAR.  Cahill--Glauber give the sharper finite-mode object:
Grassmann-valued displacement operators built from concrete CAR generators.
Their multiplier is valued in the even Grassmann algebra, not generally in the
complex phases.  Therefore this is not, without extending the scalar algebra,
the same as the ordinary projective unitary representation of the additive
group \(\operatorname{Map}(X,V)\) used in the bosonic and arithmetic Weyl
constructions.  The naive target-replacement conjecture is rejected only in
that strict scalar-projective sense.

\subsection{Current Dictionary}

The comparison is therefore compact:
\begin{itemize}
  \item finite boson: labels \(\operatorname{Map}(X,\mathbb R_q\oplus\mathbb R_p)\),
  alternating symplectic form, CCR/Weyl quantization;
  \item finite arithmetic field: labels \(E\leq\operatorname{Map}(X,V_{\mathbb F_q})\),
  alternating symplectic form after radical reduction, finite Weyl quantization;
  \item finite fermion: labels \(\operatorname{Re}(Z\oplus\overline Z)\) with
  \(Z=\operatorname{Map}(X,K)\), positive symmetric form, CAR/Clifford
  quantization.
\end{itemize}
Thus the bosonic conjecture is checked at finite kinematic level.  The
fermionic conjecture is repaired: standard fermions first enter this programme
through CAR/Clifford data on an antisymmetric Fock space, and Grassmann
variables then give the Cahill--Glauber displacement calculus in the
Grassmann-extended CAR algebra.  Section~\ref{sec:grassmann-weyl-fermion-displacements}
proves that positive
representation theorem and records the ordinary complex-Hilbert-space caveat.
